\section{Executive Summary}

\subsection{Project Overview}

The BAKO Motors Monitoring System is a full-stack embedded IoT platform developed to provide real-time diagnostics and telemetry for BAKO Motors' electric vehicle prototype. The system integrates an ESP32-based custom hardware board with a CAN bus interface, a twelve-sensor network, a live WebSocket web dashboard, and a cloud database pipeline through GSM/GPRS connectivity.

The project was designed and developed by a seven-member engineering student team at MedTech Mediterranean Institute of Technology (SMU) as part of the ISS~496 senior project (one academic semester), in direct collaboration with BAKO Motors as the industrial partner. The system is intended to serve both as a diagnostic tool for engineers and as a foundation for future predictive maintenance capabilities.

\subsection{Client Brief \& Objectives}

BAKO Motors identified several critical gaps in their existing electric vehicle prototype:

\begin{itemize}
    \item the CAN bus communication was limited to the battery management system only;
    \item no real-time data access was available to engineers, and no monitoring interface existed for vehicle diagnostics;
    \item sensors were exclusively placed on the battery pack with no coverage of the motor, MPPT controller, or power distribution unit;
    \item insufficient system documentation was available for maintenance and development teams.
\end{itemize}

The primary objectives of this project were therefore: to extend CAN bus monitoring across the full vehicle by adding external sensors; to develop a real-time web-based dashboard accessible over Wi-Fi and GPRS; to design and fabricate a custom PCB integrating all monitoring electronics; to implement a cloud telemetry pipeline for remote data storage and analysis; and to produce a Python-based software emulator enabling full system testing without physical hardware.

\subsection{Key Milestones \& Timeline}

The project was structured across a one-semester academic timeline following an iterative Agile development approach. The first phase covered system design, component selection, and feasibility study, culminating in a validated bill of materials and system architecture synoptic. The second phase focused on hardware development, firmware implementation, dashboard development, and cloud pipeline integration.

Key milestones achieved at the midway stage include: completion of the system architecture and synoptic design; full PCB schematic and layout in KiCad (Revision RECTIFICATION\_2, April~2026); implementation of CAN bus communication with the BMS via MCP2515; DS18B20 temperature sensor integration for motor, MPPT, and DC/DC converter monitoring; current and voltage measurement across all energy nodes; a Python-based sensor emulator with seven automated test scenarios; a real-time WebSocket dashboard with data export functionality; and a cloud telemetry pipeline from ESP32 through SIM800L to a SQLite database on a VPS.

All five items previously identified as remaining work at the midway stage have since been completed: the SIM800L GSM/GPRS cellular pipeline is fully operational from ESP32 to VPS; fault detection with SMS alert dispatch via threshold monitoring is implemented; the PCB was installed in the vehicle chassis and field-validated; and all testing and quality assurance procedures were completed at the Sprint~8 Production Readiness Review. GNSS position tracking integration was deprioritised following the MoSCoW review and deferred to the short-term development roadmap (Section~\ref{sec:recommendations}).

\subsection{Document Conventions}

This report follows the MedTech ISS~496 reporting standard and all electronic references use the IEEE citation style as specified in the MedTech Internship Guide. Component references throughout the document follow KiCad designator conventions, where U denotes integrated circuits, J denotes connectors, R denotes resistors, C denotes capacitors, D denotes diodes, and F denotes fuses.

Voltage levels are expressed in volts~(V) or millivolts~(mV). Temperature values are expressed in degrees Celsius~(\textdegree{}C). Current values are expressed in amperes~(A). CAN frame identifiers are expressed in hexadecimal notation with the \texttt{0x} prefix. All PCB dimensions are given in millimetres~(mm) and electrical trace widths in micrometres~(\textmu{}m) unless otherwise stated.
